% SPAR Report: Correlated Cubic Features
% Built from the SPAR report template:
%   https://github.com/the-kairos-project/spar-report-template
% Compatible with arXiv (TeX Live 2023+, pdflatex)
%
% To compile: latexmk -pdf report.tex

% ============================================================================
% DOCUMENT CLASS
% ============================================================================
% Options: midterm | final | arxiv
\documentclass[midterm]{sparreport}

% ============================================================================
% PROJECT PACKAGES AND NOTATION (matching main.tex)
% ============================================================================
\usepackage{bm}
\usepackage{xcolor}
\usepackage{enumitem}
\usepackage{subcaption}

\newtheorem{theorem}{Theorem}[section]
\newtheorem{lemma}[theorem]{Lemma}
\newtheorem{proposition}[theorem]{Proposition}
\newtheorem{corollary}[theorem]{Corollary}
\theoremstyle{definition}
\newtheorem{definition}[theorem]{Definition}
\newtheorem{example}[theorem]{Example}
\newtheorem{remark}[theorem]{Remark}

\newcommand{\R}{\mathbb{R}}
\newcommand{\T}{^\top}
\newcommand{\mc}[1]{\mathcal{#1}}
\newcommand{\E}{\mathbb{E}}
\newcommand{\V}{\mathcal{V}}
\newcommand{\Cov}{\operatorname{Cov}}
\newcommand{\Tr}{\operatorname{Tr}}
\newcommand{\Id}{\mathrm{Id}}
\newcommand{\argmin}{\operatorname*{arg\,min}}
\newcommand{\norm}[1]{\left\lVert #1 \right\rVert}
\newcommand{\inner}[2]{\left\langle #1,#2 \right\rangle}
\newcommand{\comm}{\mathcal{C}}

% ============================================================================
% METADATA
% ============================================================================
\title{Correlated Cubic Features}

\author{
  \textbf{Jeremy Atkins}\\
  \texttt{jeremy\_atkins@brown.edu}\\
  Brown University
  % Add co-authors and mentor(s) with \and:
  % \and
  % \textbf{Aniket ...}\\
  % \texttt{...}
  % \and
  % \textbf{Mentor Name}\\
  % \texttt{mentor@example.org}
}

\sparround{Fall 2026}

% ============================================================================
% DOCUMENT BEGINS
% ============================================================================
\begin{document}

\maketitle

% ============================================================================
% ABSTRACT (100-200 words)
% ============================================================================
% Overview of the research question, why it matters, approach, and findings.
% Keep it accessible to readers outside the subfield. No footnotes or citations.
\begin{abstract}
We extend Nisch's two-feature calculation to $k$ correlated cubic target
features at Bayesian equilibrium. We ask when a single committee of neurons
serves all the features, when it splits, and which populations of neurons
minimize the posterior cost after splitting.
% TODO: expand to 100-200 words; state why this matters for AI safety and
% summarize the main results.
\end{abstract}

\keywords{feature learning \and Bayesian neural networks \and superposition \and statistical mechanics}

% ============================================================================
% INTRODUCTION (100-350 words)
% ============================================================================
% The problem, why it matters, the gap it fills, and what success looks like.
\section{Introduction}

Neural networks can learn features either by aligning small numbers of highly correlated neurons (``committees''), or through small adjustments in many neurons.
The former is a form of lottery ticket hypothesis: it is possible that features and tasks are largely learned by subnetworks that are small compared to the size of the whole language model.
This would have strong implications for interpretability: namely, that we should be looking for some kind of ``core mechanism'' for a task, and that other neurons may only add minimal contribution.

We apply methods from statistical mechanics to understand the neuron structure for a two-layer network with cubic Hermite polynomial nonlinearity.
The Gaussian orthogonality of Hermite polynomials aid in solving the ground state of the system.
Based on prior work from Mainali (unpublished), we begin from a formulation of the training dynamics of the network as a free energy minimization, and then solve for the minima for cases with various kinds of symmetry.
Our primary aim is to understand how the network might perform on realistic datasets, which display multiple levels of organization.
Thus, we study the case of a teacher network which represents explicitly a lower level of organization, with the higher level represented implicitly by correlations in teacher features.
We show that the phase structure of representations learned by the student network is surprisingly complex.
We discuss implications of this result for interpretability.

% ============================================================================
% METHODOLOGY (200-500 words)
% ============================================================================
% For theoretical work: framework, key assumptions, and approach to proofs.
% Source material: main.tex, Setup through "Reduced description in terms of
% committees".
\section{Methodology}

Let $d$ be the input dimension, $K$ be the number of hidden neurons, and $P$ be the number of training examples.

Take
\begin{equation}
    K = \frac{d^2}{\Delta},\quad P = \chi K,
\end{equation}
with $\Delta, \chi > 0$ fixed.
The teacher model is defined as
\begin{align*}
    f_*(x) = \beta \sum_{j=1}^k \psi_3(e_j^\top x)
\end{align*}
where $\psi_3$ is the cubic Hermite polynomial and $\beta \geq 0$ is the output strength.
The student network is then
\begin{align*}
    f(x) = \frac{1}{\sqrt{d}} \sum_{i=1}^{K} a_i \psi_3(w_i^\top x)
\end{align*}
where the $a_i$ are the readout strength with $|a_i| \leq A$ and $N^{-1} \sum_i a_i^2 = \tau^2$, such that the readouts are fixed on the sphere $S^{K-1}$.
$w_i$ are the hidden weights for each neuron.

Let the posterior over readouts and hidden weights be
\begin{align*}
    \Pi(a, W \,|\, X, y)
    \propto \Pi_0(a, W) \,\exp\left[ - \frac{|y - f_{a, W}(X)|^2|}{2\kappa^2} \right]
\end{align*}
We can show that, for $\comm = \{n, a, r, u\}$ a proposed collection of $n$ neurons such that each weight can be written $w = ru + \sqrt{1-r^2}w_\perp$ and each have output strength $a$, the posterior for that group is
\begin{align*}
    \Pi(\comm \,|\, X, y)
    \propto \Pi_0(\comm) \, \exp\left[ -d \sum_j n_j J(r_j) -\frac{dg}{2} \E_x (f_* - f_\comm)^2 \right]
\end{align*}
where $g = \alpha/v$, and $v$ is the total variance per sample, and $J(r) = -1/2 \log(1-r^2)$.

Various manipulations can then be done to show that this can be written
\begin{align*}
    \Pi(\{b_j, u_j\} \,|\, X, y)
    \propto e^{-Ng(F - \min F)} \\
    \boxed{F = \lambda \sum_j b_j + \frac{1}{2} \E_x (f_* - f_\comm)^2}
\end{align*}
where the boxed expression is the \emph{free energy}.
$\lambda$ is an $O(1)$ positive constant derived from other values, and the $b_j = O(n d^{-1/2})$ are strengths.
The training process will tend to minimize the free energy.

% ============================================================================
% PRELIMINARY RESULTS (200-500 words)
% ============================================================================
% Key insights and partial results.
% Source material: main.tex ("Three Correlated Features", "k Features",
% "Local Stability with Constant Row Sums", "Exact Splitting of Orthogonal
% Feature Groups"), isosceles_setup.tex, orthogonal_extension.tex.
\section{Preliminary Results}

% \subsection{Critical correlation for $k$ symmetric features}
% \subsection{Stability of the shared direction}
% \subsection{Exact splitting of orthogonal feature groups}

% ============================================================================
% CHALLENGES & OPEN QUESTIONS (100-300 words, optional but encouraged)
% ============================================================================
\section{Challenges \& Open Questions}

% TODO

% ============================================================================
% NEXT STEPS (150-400 words)
% ============================================================================
\section{Next Steps}

% TODO

% ============================================================================
% BIBLIOGRAPHY
% ============================================================================
% For arXiv submission, include the generated .bbl file.
% Uncomment once references.bib has entries and the text cites them
% (bibtex errors out on a document with no citations).
% \bibliographystyle{plain}
% \bibliography{references}

\end{document}
